\subsection{Materia necesaria}

\begin{definicion}{Capacidad con merma y mezcla}
En procesos con bifurcaciones, cada flujo debe seguir su rendimiento. Si \(F\) kg/h entran a secado y luego \(80\%\) va a molido y \(20\%\) a azucarado con \(10\%\) de descarte, el flujo efectivo que llega a mezcla es:
\[
0.8F+0.2F(0.9)=0.98F.
\]
\end{definicion}

\begin{formulas}{Revision periodica y continua}
Revision periodica con periodo \(T\) y lead time \(L\):
\[
S_T=d(T+L)+z\sigma\sqrt{T+L}.
\]
Revision continua:
\[
Q^*=\sqrt{\frac{2DS}{h}},
\qquad
ROP=dL+z\sigma\sqrt L.
\]
\end{formulas}

\begin{formulas}{Wagner-Whitin}
Si se ordena en \(i\) para cubrir hasta \(j\):
\[
C(i,j)=S+h\sum_{k=i}^{j}(k-i)D_k.
\]
La recurrencia es:
\[
F(j)=\min_{1\le i\le j}\{F(i-1)+C(i,j)\}.
\]
\end{formulas}

\begin{algoritmo}{Capacidad e inventarios}
\begin{enumerate}
    \item Definir variables de flujo antes de calcular capacidades.
    \item Escribir restricciones de cada maquina con el flujo que realmente procesa.
    \item Identificar cuello de botella por maximizacion, no por inspeccion visual.
    \item En inventarios, fijar unidad temporal unica antes de usar formulas.
\end{enumerate}
\end{algoritmo}
